% --- Archon physics formalization source begin ---
% archon:physics
% archon:covers PhyXMiniProblems/problem_phyx_mini_0936.lean
% archon:source-report reports/phyx_mini/problem_phyx_mini_0936.source.json

\chapter{Physics problem phyx\_mini\_0936}
\label{ch:PhyXMiniProblems_problem_phyx_mini_0936}

\paragraph{Problem source.}
\#\# Physical scenario

A 500-loop circular wire coil with radius \textbackslash{}( 4.00 \textbackslash{} \textbackslash{}mathrm\{cm\} \textbackslash{}) is placed between the poles of a large electromagnet. The magnetic field is uniform and makes an angle of \textbackslash{}( 60\textasciicircum{}\textbackslash{}circ \textbackslash{}) with the plane of the coil; it decreases at \textbackslash{}( 0.200 \textbackslash{} \textbackslash{}mathrm\{T/s\} \textbackslash{}).

\#\# Question

What are the magnitude and direction of the induced emf?

\#\# Answer choices

- A: A: 0.456V

- B: B: 0.435V

- C: C: 0.544V

- D: D: 0.231V

\#\# Figure caption (auxiliary; use the image as primary evidence)

The image is a diagram depicting a magnetic field and a current-carrying wire. Here’s a detailed breakdown:

- **Magnetic Field (\textbackslash{}( \textbackslash{}vec\{B\} \textbackslash{})):** It is represented by five horizontal lines with arrows pointing from left to right, indicating the direction of the field from the North (N) side to the South (S) side.

- **Wire:** A diagonal wire is shown with a circular cross-section at its top end, indicating current flow out of the plane. The wire is labeled with a length of "4.00 cm".

- **Angles:** 
  - \textbackslash{}(60\textasciicircum{}\textbackslash{}circ\textbackslash{}) is the angle between the horizontal magnetic field lines and the wire.
  - \textbackslash{}( \textbackslash{}varnothing = 30\textasciicircum{}\textbackslash{}circ \textbackslash{}) is the angle between the wire and a vector \textbackslash{}( \textbackslash{}vec\{A\} \textbackslash{}) parallel to the magnetic field lines.

- **Boundaries:** The left and right sides of the diagram are labeled as N (North) and S (South), respectively, enclosing the magnetic field region.

The diagram illustrates the orientation of a current-carrying wire in a magnetic field and uses vectors and angles to indicate directions and relationships between elements.

\#\# Dataset metadata

Electromagnetic Induction; Physical Model Grounding Reasoning, Multi-Formula Reasoning

\paragraph{Recorded answer/context.}
B: B: 0.435V

\paragraph{Figure/image path.}
/root/proposal\_for\_physic/science-mango/phyx\_mini\_run/phyx\_data/test\_image/936.png

\paragraph{Formalization target.}
create a compiling Lean file with sorry bodies at `PhyXMiniProblems/problem\_phyx\_mini\_0936.lean`. The Lean declarations must preserve the physical quantities, dimensions or dimensional roles, figure labels, governing-law hypotheses, and final relation expressed by this problem.
Use Mathlib/Physlib names found through LeanExplore where available. If a physics API is missing, introduce faithful local abstractions rather than scalar placeholder aliases.

\begin{theorem}[Physics formalization target]
\label{thm:physics:phyx_mini_0936:target}
\lean{PhyXMiniProblems.ProblemPhyXMini0936.problem_phyx_mini_0936}
\uses{def:physics:phyx-mini-0936:phyxminiproblems-problemphyxmini0936-signedemfinvolts, def:physics:phyx-mini-0936:phyxminiproblems-problemphyxmini0936-circularcoilinductionsetup, def:physics:phyx-mini-0936:phyxminiproblems-problemphyxmini0936-matchesprimarycircularcoilfigure, def:physics:phyx-mini-0936:phyxminiproblems-problemphyxmini0936-matchescircularcoilproblemdata, def:physics:phyx-mini-0936:phyxminiproblems-problemphyxmini0936-hasphysicalcircularcoilparameters, def:physics:phyx-mini-0936:phyxminiproblems-problemphyxmini0936-hasuniformappliedmagneticfield, def:physics:phyx-mini-0936:phyxminiproblems-problemphyxmini0936-satisfiescircularcoilarealaw, def:physics:phyx-mini-0936:phyxminiproblems-problemphyxmini0936-satisfiesuniformfieldfluxlaws, def:physics:phyx-mini-0936:phyxminiproblems-problemphyxmini0936-satisfiesfaradaylenzlaw, def:physics:phyx-mini-0936:phyxminiproblems-problemphyxmini0936-recordeddatasetanswer, def:physics:phyx-mini-0936:phyxminiproblems-problemphyxmini0936-isuniquenearestdisplayedchoice}
The assigned autoformalize agent should translate this physics problem into Lean declarations in the covered file, with theorem and lemma proof bodies written as `by sorry`.
\end{theorem}
\begin{proof}
This is an autoformalization task, not a proof task. Produce statements that can later be proved by the physics prover without weakening the physical model.
\end{proof}

% --- Archon declaration topology begin ---
\section{Declaration topology}
\begin{definition}[magnetic Flux Density Dimension]
  \label{def:physics:phyx-mini-0936:phyxminiproblems-problemphyxmini0936-magneticfluxdensitydimension}
  \lean{PhyXMiniProblems.ProblemPhyXMini0936.magneticFluxDensityDimension}
  The dimension M T\textasciicircum{}-1 C\textasciicircum{}-1 of magnetic flux density (tesla)
\end{definition}
\begin{proof}
  This is the declared model definition of magnetic Flux Density Dimension.
\end{proof}
\begin{definition}[magnetic Flux Density Rate Dimension]
  \label{def:physics:phyx-mini-0936:phyxminiproblems-problemphyxmini0936-magneticfluxdensityratedimension}
  \lean{PhyXMiniProblems.ProblemPhyXMini0936.magneticFluxDensityRateDimension}
  The dimension M T\textasciicircum{}-2 C\textasciicircum{}-1 of a magnetic-flux-density rate
\end{definition}
\begin{proof}
  This is the declared model definition of magnetic Flux Density Rate Dimension.
\end{proof}
\begin{definition}[magnetic Flux Dimension]
  \label{def:physics:phyx-mini-0936:phyxminiproblems-problemphyxmini0936-magneticfluxdimension}
  \lean{PhyXMiniProblems.ProblemPhyXMini0936.magneticFluxDimension}
  The dimension M L\textasciicircum{}2 T\textasciicircum{}-1 C\textasciicircum{}-1 of magnetic flux (weber)
\end{definition}
\begin{proof}
  This is the declared model definition of magnetic Flux Dimension.
\end{proof}
\begin{definition}[electromotive Force Dimension]
  \label{def:physics:phyx-mini-0936:phyxminiproblems-problemphyxmini0936-electromotiveforcedimension}
  \lean{PhyXMiniProblems.ProblemPhyXMini0936.electromotiveForceDimension}
  The dimension M L\textasciicircum{}2 T\textasciicircum{}-2 C\textasciicircum{}-1 shared by flux rate and emf (volt)
\end{definition}
\begin{proof}
  This is the declared model definition of electromotive Force Dimension.
\end{proof}
\begin{definition}[Length Magnitude]
  \label{def:physics:phyx-mini-0936:phyxminiproblems-problemphyxmini0936-lengthmagnitude}
  \lean{PhyXMiniProblems.ProblemPhyXMini0936.LengthMagnitude}
  A nonnegative, unit-independent physical length
\end{definition}
\begin{proof}
  This is the declared model definition of Length Magnitude.
\end{proof}
\begin{definition}[Area Magnitude]
  \label{def:physics:phyx-mini-0936:phyxminiproblems-problemphyxmini0936-areamagnitude}
  \lean{PhyXMiniProblems.ProblemPhyXMini0936.AreaMagnitude}
  A nonnegative, unit-independent physical area
\end{definition}
\begin{proof}
  This is the declared model definition of Area Magnitude.
\end{proof}
\begin{definition}[Magnetic Flux Density Magnitude]
  \label{def:physics:phyx-mini-0936:phyxminiproblems-problemphyxmini0936-magneticfluxdensitymagnitude}
  \lean{PhyXMiniProblems.ProblemPhyXMini0936.MagneticFluxDensityMagnitude}
  \uses{def:physics:phyx-mini-0936:phyxminiproblems-problemphyxmini0936-magneticfluxdensitydimension}
  A nonnegative, unit-independent magnetic-flux-density magnitude
\end{definition}
\begin{proof}
  This is the declared model definition of Magnetic Flux Density Magnitude.
\end{proof}
\begin{definition}[Signed Magnetic Flux Density Rate]
  \label{def:physics:phyx-mini-0936:phyxminiproblems-problemphyxmini0936-signedmagneticfluxdensityrate}
  \lean{PhyXMiniProblems.ProblemPhyXMini0936.SignedMagneticFluxDensityRate}
  \uses{def:physics:phyx-mini-0936:phyxminiproblems-problemphyxmini0936-magneticfluxdensityratedimension}
  A signed, unit-independent rate of change of magnetic flux density
\end{definition}
\begin{proof}
  This is the declared model definition of Signed Magnetic Flux Density Rate.
\end{proof}
\begin{definition}[Signed Magnetic Flux]
  \label{def:physics:phyx-mini-0936:phyxminiproblems-problemphyxmini0936-signedmagneticflux}
  \lean{PhyXMiniProblems.ProblemPhyXMini0936.SignedMagneticFlux}
  \uses{def:physics:phyx-mini-0936:phyxminiproblems-problemphyxmini0936-magneticfluxdimension}
  A signed magnetic flux through one oriented turn of the coil
\end{definition}
\begin{proof}
  This is the declared model definition of Signed Magnetic Flux.
\end{proof}
\begin{definition}[Signed Magnetic Flux Rate]
  \label{def:physics:phyx-mini-0936:phyxminiproblems-problemphyxmini0936-signedmagneticfluxrate}
  \lean{PhyXMiniProblems.ProblemPhyXMini0936.SignedMagneticFluxRate}
  \uses{def:physics:phyx-mini-0936:phyxminiproblems-problemphyxmini0936-electromotiveforcedimension}
  A signed rate of magnetic flux through one oriented turn
\end{definition}
\begin{proof}
  This is the declared model definition of Signed Magnetic Flux Rate.
\end{proof}
\begin{definition}[Signed Electromotive Force]
  \label{def:physics:phyx-mini-0936:phyxminiproblems-problemphyxmini0936-signedelectromotiveforce}
  \lean{PhyXMiniProblems.ProblemPhyXMini0936.SignedElectromotiveForce}
  \uses{def:physics:phyx-mini-0936:phyxminiproblems-problemphyxmini0936-electromotiveforcedimension}
  A signed induced electromotive force around the oriented winding
\end{definition}
\begin{proof}
  This is the declared model definition of Signed Electromotive Force.
\end{proof}
\begin{definition}[length In Meters]
  \label{def:physics:phyx-mini-0936:phyxminiproblems-problemphyxmini0936-lengthinmeters}
  \lean{PhyXMiniProblems.ProblemPhyXMini0936.lengthInMeters}
  \uses{def:physics:phyx-mini-0936:phyxminiproblems-problemphyxmini0936-lengthmagnitude}
  Read a physical length in coherent-SI metres
\end{definition}
\begin{proof}
  This is the declared model definition of length In Meters.
\end{proof}
\begin{definition}[length In Centimeters]
  \label{def:physics:phyx-mini-0936:phyxminiproblems-problemphyxmini0936-lengthincentimeters}
  \lean{PhyXMiniProblems.ProblemPhyXMini0936.lengthInCentimeters}
  \uses{def:physics:phyx-mini-0936:phyxminiproblems-problemphyxmini0936-lengthmagnitude, def:physics:phyx-mini-0936:phyxminiproblems-problemphyxmini0936-lengthinmeters}
  Read a physical length in centimetres
\end{definition}
\begin{proof}
  This is the declared model definition of length In Centimeters.
\end{proof}
\begin{definition}[area In Square Meters]
  \label{def:physics:phyx-mini-0936:phyxminiproblems-problemphyxmini0936-areainsquaremeters}
  \lean{PhyXMiniProblems.ProblemPhyXMini0936.areaInSquareMeters}
  \uses{def:physics:phyx-mini-0936:phyxminiproblems-problemphyxmini0936-areamagnitude}
  Read a physical area in coherent-SI square metres
\end{definition}
\begin{proof}
  This is the declared model definition of area In Square Meters.
\end{proof}
\begin{definition}[magnetic Flux Density In Teslas]
  \label{def:physics:phyx-mini-0936:phyxminiproblems-problemphyxmini0936-magneticfluxdensityinteslas}
  \lean{PhyXMiniProblems.ProblemPhyXMini0936.magneticFluxDensityInTeslas}
  \uses{def:physics:phyx-mini-0936:phyxminiproblems-problemphyxmini0936-magneticfluxdensitymagnitude}
  Read magnetic flux density in coherent-SI teslas
\end{definition}
\begin{proof}
  This is the declared model definition of magnetic Flux Density In Teslas.
\end{proof}
\begin{definition}[magnetic Flux Density Rate In Teslas Per Second]
  \label{def:physics:phyx-mini-0936:phyxminiproblems-problemphyxmini0936-magneticfluxdensityrateinteslaspersecond}
  \lean{PhyXMiniProblems.ProblemPhyXMini0936.magneticFluxDensityRateInTeslasPerSecond}
  \uses{def:physics:phyx-mini-0936:phyxminiproblems-problemphyxmini0936-signedmagneticfluxdensityrate}
  Read a signed magnetic-flux-density rate in teslas per second
\end{definition}
\begin{proof}
  This is the declared model definition of magnetic Flux Density Rate In Teslas Per Second.
\end{proof}
\begin{definition}[signed Magnetic Flux In Webers]
  \label{def:physics:phyx-mini-0936:phyxminiproblems-problemphyxmini0936-signedmagneticfluxinwebers}
  \lean{PhyXMiniProblems.ProblemPhyXMini0936.signedMagneticFluxInWebers}
  \uses{def:physics:phyx-mini-0936:phyxminiproblems-problemphyxmini0936-signedmagneticflux}
  Read signed magnetic flux in coherent-SI webers
\end{definition}
\begin{proof}
  This is the declared model definition of signed Magnetic Flux In Webers.
\end{proof}
\begin{definition}[signed Magnetic Flux Rate In Webers Per Second]
  \label{def:physics:phyx-mini-0936:phyxminiproblems-problemphyxmini0936-signedmagneticfluxrateinweberspersecond}
  \lean{PhyXMiniProblems.ProblemPhyXMini0936.signedMagneticFluxRateInWebersPerSecond}
  \uses{def:physics:phyx-mini-0936:phyxminiproblems-problemphyxmini0936-signedmagneticfluxrate}
  Read a signed magnetic-flux rate in webers per second
\end{definition}
\begin{proof}
  This is the declared model definition of signed Magnetic Flux Rate In Webers Per Second.
\end{proof}
\begin{definition}[signed Emf In Volts]
  \label{def:physics:phyx-mini-0936:phyxminiproblems-problemphyxmini0936-signedemfinvolts}
  \lean{PhyXMiniProblems.ProblemPhyXMini0936.signedEmfInVolts}
  \uses{def:physics:phyx-mini-0936:phyxminiproblems-problemphyxmini0936-signedelectromotiveforce}
  Read a signed induced emf in coherent-SI volts
\end{definition}
\begin{proof}
  This is the declared model definition of signed Emf In Volts.
\end{proof}
\begin{definition}[Magnet Pole]
  \label{def:physics:phyx-mini-0936:phyxminiproblems-problemphyxmini0936-magnetpole}
  \lean{PhyXMiniProblems.ProblemPhyXMini0936.MagnetPole}
  The two labeled pole pieces in the supplied image
\end{definition}
\begin{proof}
  This is the declared model definition of Magnet Pole.
\end{proof}
\begin{definition}[Horizontal Side]
  \label{def:physics:phyx-mini-0936:phyxminiproblems-problemphyxmini0936-horizontalside}
  \lean{PhyXMiniProblems.ProblemPhyXMini0936.HorizontalSide}
  Left and right sides of the primary image
\end{definition}
\begin{proof}
  This is the declared model definition of Horizontal Side.
\end{proof}
\begin{definition}[Horizontal Direction]
  \label{def:physics:phyx-mini-0936:phyxminiproblems-problemphyxmini0936-horizontaldirection}
  \lean{PhyXMiniProblems.ProblemPhyXMini0936.HorizontalDirection}
  Horizontal arrow directions used by the field lines
\end{definition}
\begin{proof}
  This is the declared model definition of Horizontal Direction.
\end{proof}
\begin{definition}[Coil Traversal Sense]
  \label{def:physics:phyx-mini-0936:phyxminiproblems-problemphyxmini0936-coiltraversalsense}
  \lean{PhyXMiniProblems.ProblemPhyXMini0936.CoilTraversalSense}
  An oriented traversal of the coil. The viewpoint is from the tip of the positive area-normal vector looking back toward the coil
\end{definition}
\begin{proof}
  This is the declared model definition of Coil Traversal Sense.
\end{proof}
\begin{definition}[orientation Sign]
  \label{def:physics:phyx-mini-0936:phyxminiproblems-problemphyxmini0936-coiltraversalsense-orientationsign}
  \lean{PhyXMiniProblems.ProblemPhyXMini0936.CoilTraversalSense.orientationSign}
  \uses{def:physics:phyx-mini-0936:phyxminiproblems-problemphyxmini0936-coiltraversalsense}
  The sign attached to an oriented traversal. This is an orientation convention, not a prediction of the requested direction
\end{definition}
\begin{proof}
  This is the declared model definition of orientation Sign.
\end{proof}
\begin{definition}[Circular Coil Figure]
  \label{def:physics:phyx-mini-0936:phyxminiproblems-problemphyxmini0936-circularcoilfigure}
  \lean{PhyXMiniProblems.ProblemPhyXMini0936.CircularCoilFigure}
  \uses{def:physics:phyx-mini-0936:phyxminiproblems-problemphyxmini0936-magnetpole, def:physics:phyx-mini-0936:phyxminiproblems-problemphyxmini0936-horizontalside, def:physics:phyx-mini-0936:phyxminiproblems-problemphyxmini0936-horizontaldirection}
  Literal presentation data transcribed from image 936.png. The diagonal segment is the edge-on depiction of the circular coil plane, with hatched endcaps; the arrow labeled A is its positive area normal
\end{definition}
\begin{proof}
  This is the declared model definition of Circular Coil Figure.
\end{proof}
\begin{definition}[Circular Coil Induction Setup]
  \label{def:physics:phyx-mini-0936:phyxminiproblems-problemphyxmini0936-circularcoilinductionsetup}
  \lean{PhyXMiniProblems.ProblemPhyXMini0936.CircularCoilInductionSetup}
  \uses{def:physics:phyx-mini-0936:phyxminiproblems-problemphyxmini0936-lengthmagnitude, def:physics:phyx-mini-0936:phyxminiproblems-problemphyxmini0936-areamagnitude, def:physics:phyx-mini-0936:phyxminiproblems-problemphyxmini0936-magneticfluxdensitymagnitude, def:physics:phyx-mini-0936:phyxminiproblems-problemphyxmini0936-signedmagneticfluxdensityrate, def:physics:phyx-mini-0936:phyxminiproblems-problemphyxmini0936-signedmagneticflux, def:physics:phyx-mini-0936:phyxminiproblems-problemphyxmini0936-signedmagneticfluxrate, def:physics:phyx-mini-0936:phyxminiproblems-problemphyxmini0936-signedelectromotiveforce, def:physics:phyx-mini-0936:phyxminiproblems-problemphyxmini0936-coiltraversalsense, def:physics:phyx-mini-0936:phyxminiproblems-problemphyxmini0936-circularcoilfigure}
  Independent physical quantities for the experiment at the instant in question. Neither the signed emf nor its traversal sense is defined from an answer choice or from the requested final relation
\end{definition}
\begin{proof}
  This is the declared model definition of Circular Coil Induction Setup.
\end{proof}
\begin{definition}[Matches Primary Circular Coil Figure]
  \label{def:physics:phyx-mini-0936:phyxminiproblems-problemphyxmini0936-matchesprimarycircularcoilfigure}
  \lean{PhyXMiniProblems.ProblemPhyXMini0936.MatchesPrimaryCircularCoilFigure}
  \uses{def:physics:phyx-mini-0936:phyxminiproblems-problemphyxmini0936-circularcoilinductionsetup}
  Primary-image evidence. These fields contain only literal labels and qualitative directions visible in the raster
\end{definition}
\begin{proof}
  This is the declared model definition of Matches Primary Circular Coil Figure.
\end{proof}
\begin{definition}[Matches Circular Coil Problem Data]
  \label{def:physics:phyx-mini-0936:phyxminiproblems-problemphyxmini0936-matchescircularcoilproblemdata}
  \lean{PhyXMiniProblems.ProblemPhyXMini0936.MatchesCircularCoilProblemData}
  \uses{def:physics:phyx-mini-0936:phyxminiproblems-problemphyxmini0936-lengthincentimeters, def:physics:phyx-mini-0936:phyxminiproblems-problemphyxmini0936-magneticfluxdensityrateinteslaspersecond, def:physics:phyx-mini-0936:phyxminiproblems-problemphyxmini0936-circularcoilinductionsetup}
  Numerical and geometric data stated in the prose and tied to the physical setup. The negative SI rate encodes the word "decreases"
\end{definition}
\begin{proof}
  This is the declared model definition of Matches Circular Coil Problem Data.
\end{proof}
\begin{definition}[Has Physical Circular Coil Parameters]
  \label{def:physics:phyx-mini-0936:phyxminiproblems-problemphyxmini0936-hasphysicalcircularcoilparameters}
  \lean{PhyXMiniProblems.ProblemPhyXMini0936.HasPhysicalCircularCoilParameters}
  \uses{def:physics:phyx-mini-0936:phyxminiproblems-problemphyxmini0936-lengthinmeters, def:physics:phyx-mini-0936:phyxminiproblems-problemphyxmini0936-areainsquaremeters, def:physics:phyx-mini-0936:phyxminiproblems-problemphyxmini0936-magneticfluxdensityinteslas, def:physics:phyx-mini-0936:phyxminiproblems-problemphyxmini0936-circularcoilinductionsetup}
  Positivity and non-vacuity conditions for the physical apparatus
\end{definition}
\begin{proof}
  This is the declared model definition of Has Physical Circular Coil Parameters.
\end{proof}
\begin{definition}[Has Uniform Applied Magnetic Field]
  \label{def:physics:phyx-mini-0936:phyxminiproblems-problemphyxmini0936-hasuniformappliedmagneticfield}
  \lean{PhyXMiniProblems.ProblemPhyXMini0936.HasUniformAppliedMagneticField}
  \uses{def:physics:phyx-mini-0936:phyxminiproblems-problemphyxmini0936-magneticfluxdensityinteslas, def:physics:phyx-mini-0936:phyxminiproblems-problemphyxmini0936-circularcoilinductionsetup}
  The Physlib vector field has the stated scalar magnitude everywhere in the uniform region at the observation time
\end{definition}
\begin{proof}
  This is the declared model definition of Has Uniform Applied Magnetic Field.
\end{proof}
\begin{definition}[Satisfies Circular Coil Area Law]
  \label{def:physics:phyx-mini-0936:phyxminiproblems-problemphyxmini0936-satisfiescircularcoilarealaw}
  \lean{PhyXMiniProblems.ProblemPhyXMini0936.SatisfiesCircularCoilAreaLaw}
  \uses{def:physics:phyx-mini-0936:phyxminiproblems-problemphyxmini0936-lengthinmeters, def:physics:phyx-mini-0936:phyxminiproblems-problemphyxmini0936-areainsquaremeters, def:physics:phyx-mini-0936:phyxminiproblems-problemphyxmini0936-circularcoilinductionsetup}
  The geometric area of a circular coil is pi r\textasciicircum{}2
\end{definition}
\begin{proof}
  This is the declared model definition of Satisfies Circular Coil Area Law.
\end{proof}
\begin{definition}[Satisfies Uniform Field Flux Laws]
  \label{def:physics:phyx-mini-0936:phyxminiproblems-problemphyxmini0936-satisfiesuniformfieldfluxlaws}
  \lean{PhyXMiniProblems.ProblemPhyXMini0936.SatisfiesUniformFieldFluxLaws}
  \uses{def:physics:phyx-mini-0936:phyxminiproblems-problemphyxmini0936-areainsquaremeters, def:physics:phyx-mini-0936:phyxminiproblems-problemphyxmini0936-magneticfluxdensityinteslas, def:physics:phyx-mini-0936:phyxminiproblems-problemphyxmini0936-magneticfluxdensityrateinteslaspersecond, def:physics:phyx-mini-0936:phyxminiproblems-problemphyxmini0936-signedmagneticfluxinwebers, def:physics:phyx-mini-0936:phyxminiproblems-problemphyxmini0936-signedmagneticfluxrateinweberspersecond, def:physics:phyx-mini-0936:phyxminiproblems-problemphyxmini0936-circularcoilinductionsetup}
  For a uniform field at fixed coil orientation, the signed flux through one turn is B A cos(theta) and its instantaneous rate is (dB/dt) A cos(theta), where theta is measured from the positive area normal. These are general flux laws and contain no final emf value
\end{definition}
\begin{proof}
  This is the declared model definition of Satisfies Uniform Field Flux Laws.
\end{proof}
\begin{definition}[Satisfies Faraday Lenz Law]
  \label{def:physics:phyx-mini-0936:phyxminiproblems-problemphyxmini0936-satisfiesfaradaylenzlaw}
  \lean{PhyXMiniProblems.ProblemPhyXMini0936.SatisfiesFaradayLenzLaw}
  \uses{def:physics:phyx-mini-0936:phyxminiproblems-problemphyxmini0936-signedmagneticfluxrateinweberspersecond, def:physics:phyx-mini-0936:phyxminiproblems-problemphyxmini0936-signedemfinvolts, def:physics:phyx-mini-0936:phyxminiproblems-problemphyxmini0936-circularcoilinductionsetup}
  Faraday's law relates the independent signed emf to the per-turn flux rate. The second field is the generic orientation-sign convention for the emf; it does not select either traversal direction in advance
\end{definition}
\begin{proof}
  This is the declared model definition of Satisfies Faraday Lenz Law.
\end{proof}
\begin{definition}[Answer Choice]
  \label{def:physics:phyx-mini-0936:phyxminiproblems-problemphyxmini0936-answerchoice}
  \lean{PhyXMiniProblems.ProblemPhyXMini0936.AnswerChoice}
  Labels of the four displayed voltage choices
\end{definition}
\begin{proof}
  This is the declared model definition of Answer Choice.
\end{proof}
\begin{definition}[emf Magnitude In Volts]
  \label{def:physics:phyx-mini-0936:phyxminiproblems-problemphyxmini0936-answerchoice-emfmagnitudeinvolts}
  \lean{PhyXMiniProblems.ProblemPhyXMini0936.AnswerChoice.emfMagnitudeInVolts}
  \uses{def:physics:phyx-mini-0936:phyxminiproblems-problemphyxmini0936-answerchoice}
  Voltage magnitude printed beside each displayed answer label
\end{definition}
\begin{proof}
  This is the declared model definition of emf Magnitude In Volts.
\end{proof}
\begin{definition}[recorded Dataset Answer]
  \label{def:physics:phyx-mini-0936:phyxminiproblems-problemphyxmini0936-recordeddatasetanswer}
  \lean{PhyXMiniProblems.ProblemPhyXMini0936.recordedDatasetAnswer}
  \uses{def:physics:phyx-mini-0936:phyxminiproblems-problemphyxmini0936-answerchoice}
  The answer label recorded in the supplied dataset metadata
\end{definition}
\begin{proof}
  This is the declared model definition of recorded Dataset Answer.
\end{proof}
\begin{definition}[Is Unique Nearest Displayed Choice]
  \label{def:physics:phyx-mini-0936:phyxminiproblems-problemphyxmini0936-isuniquenearestdisplayedchoice}
  \lean{PhyXMiniProblems.ProblemPhyXMini0936.IsUniqueNearestDisplayedChoice}
  \uses{def:physics:phyx-mini-0936:phyxminiproblems-problemphyxmini0936-answerchoice}
  A choice is uniquely nearest to a computed magnitude when it has strictly smaller absolute error than every differently labeled choice
\end{definition}
\begin{proof}
  This is the declared model definition of Is Unique Nearest Displayed Choice.
\end{proof}
% --- Archon declaration topology end ---
% --- Archon physics formalization source end ---
